\documentclass[11pt,a4paper,oneside]{book}

\usepackage{fontspec}
\usepackage{polyglossia}
\setmainlanguage{russian}
\setotherlanguage{english}
\setmainfont{DejaVu Serif}
\setsansfont{DejaVu Sans}
\setmonofont{DejaVu Sans Mono}

\usepackage[a4paper,margin=24mm]{geometry}
\usepackage{amsmath,amssymb}
\usepackage{array,booktabs,longtable}
\usepackage{xcolor}
\usepackage{hyperref}
\usepackage{listings}
\usepackage{enumitem}

\hypersetup{
  colorlinks=true,
  linkcolor=blue!45!black,
  urlcolor=blue!45!black,
  pdftitle={ADDA Neural Preconditioner: техническая книга проекта},
  pdfauthor={Project notes}
}

\lstdefinestyle{code}{
  basicstyle=\ttfamily\small,
  breaklines=true,
  columns=fullflexible,
  keepspaces=true,
  frame=single,
  framerule=0.3pt,
  rulecolor=\color{black!25},
  backgroundcolor=\color{black!3},
  xleftmargin=1em,
  xrightmargin=1em
}
\lstset{style=code}

\newcommand{\code}[1]{\texttt{\detokenize{#1}}}
\newcommand{\file}[1]{\texttt{\detokenize{#1}}}
\newcommand{\term}[1]{\textbf{#1}}
\newcommand{\R}{\mathbb{R}}
\newcommand{\C}{\mathbb{C}}

\title{
  \Huge ADDA Neural Preconditioner\\
  \Large техническая книга проекта
}
\author{Состояние рабочей копии: \code{convsai-universal}}
\date{1 июня 2026}

\begin{document}
\frontmatter
\maketitle
\tableofcontents

\chapter*{Зачем нужен этот документ}
\addcontentsline{toc}{chapter}{Зачем нужен этот документ}

В этом проекте смешаны три уровня сложности: численный метод DDA, модификация
ADDA на C и нейросетевое обучение предобуславливателя. Если смотреть только на
логи обучения, легко сделать неверный вывод: \emph{loss стал лучше, значит
ADDA ускорится}. На практике это не так. Модель может уменьшать внутренний
training loss и одновременно делать реальный расчёт медленнее.

Цель этой книги --- дать не поверхностную карту файлов, а рабочую модель
понимания:

\begin{itemize}[leftmargin=2em]
  \item что именно решает ADDA;
  \item что такое предобуславливатель и почему он может ускорять или замедлять;
  \item чем отличаются ConvSAI/K2 и Spectral ConvSAI;
  \item как устроены входы, выходы, loss-функции и экспорт;
  \item как читать логи, heatmap и wall-time метрики;
  \item почему последние ConvSAI/K2 fine-tune попытки были подозрительными;
  \item какой порядок действий нужен, чтобы честно доказать полезность модели.
\end{itemize}

\mainmatter

\chapter{Короткая картина проекта}

\section{Что мы ускоряем}

ADDA решает систему линейных уравнений, возникающую в DDA
(\foreignlanguage{english}{Discrete Dipole Approximation}):

\begin{equation}
  A x = b,
\end{equation}

где $A$ --- большая комплексная матрица взаимодействия диполей, $x$ ---
неизвестные дипольные поляризации, $b$ --- правая часть, задаваемая падающей
электромагнитной волной. Матрица не формируется как плотная матрица; её
применение реализуется через FFT.

Главная цена расчёта --- итерационный решатель. В проекте чаще всего
используется \code{BiCGStab}, а для некоторых больших preconditioner-heavy
случаев тестируется \code{FGMRES}. Если без предобуславливателя нужно сотни
или тысячи итераций, хороший предобуславливатель может уменьшить это число в
разы или десятки раз.

\section{Что такое нейросетевой предобуславливатель}

Предобуславливатель --- оператор $M$, который должен примерно вести себя как
$A^{-1}$. Вместо исходной системы решатель работает с более удобной системой:

\begin{equation}
  M A x = M b.
\end{equation}

Если $M A \approx I$, спектр системы становится лучше, и Krylov-решатель
сходится быстрее. Но $M$ сам стоит денег: его надо импортировать, хранить и
применять на каждом шаге. Поэтому реальная цель:

\begin{equation}
  T_{\mathrm{precond}} < T_{\mathrm{baseline}},
\end{equation}

а не просто меньшее число итераций или красивый loss.

\section{Две ветки}

Проект разделён на две независимые ветки.

\begin{description}[leftmargin=3.5cm,style=nextline]
  \item[\code{convsai-universal}]
  Текущая рабочая копия. Содержит ConvSAI Universal, также называемый K2 или
  K$^2$ v3. Это пространственный свёрточный предобуславливатель.

  \item[\code{spectral-convsai}]
  Отдельная ветка со Spectral ConvSAI. Она предсказывает предобуславливатель
  прямо в частотной области и лучше подходит для больших grid, но требует
  problem-specific export.
\end{description}

\chapter{Численная основа}

\section{DDA и матричный оператор}

В DDA частица заменяется набором диполей. Для каждого диполя есть три
компоненты поляризации, поэтому размер системы равен:

\begin{equation}
  n = 3N,
\end{equation}

где $N$ --- число диполей. Оператор обычно можно представить как:

\begin{equation}
  A = I - \alpha G,
\end{equation}

где $\alpha$ --- поляризуемость, а $G$ --- оператор взаимодействия диполей.
В коде Python эта операция реализована в \file{core/fft_matvec.py}. В ADDA
она реализована в нативном C-коде.

\section{Почему FFT важен}

Наивное применение плотной матрицы $A$ стоило бы $O(N^2)$. За счёт
свёрточной структуры взаимодействия применение $A v$ можно делать через FFT:

\begin{equation}
  A v = v - \alpha \mathcal{F}^{-1}\!\left(\widehat{G}\,\mathcal{F}(v)\right).
\end{equation}

Это делает один matvec намного дешевле, но всё равно при сотнях итераций он
остаётся главным временем расчёта.

\section{Параметры grid, dpl, kd и m}

\begin{longtable}{p{0.22\textwidth}p{0.68\textwidth}}
\toprule
Параметр & Значение в проекте \\
\midrule
\code{grid} & Размер дискретизации формы. Например \code{grid=80}. \\
\code{dpl} & Dipoles per wavelength в ADDA. Часто используется \code{-dpl 15}. \\
\code{kd} & Параметр обучения Python. В этом проекте \code{kd = 2*pi/dpl}. \\
\code{m} & Комплексный показатель преломления: $m=m_{\mathrm{re}}+i m_{\mathrm{im}}$. \\
\code{eps} & Точность ADDA. \code{-eps 3} соответствует порядку $10^{-3}$. \\
\bottomrule
\end{longtable}

Для \code{dpl=15}:

\begin{equation}
  kd = \frac{2\pi}{15} = 0.41887902047863906.
\end{equation}

Экспорт нейросетевого предобуславливателя и запуск ADDA должны использовать
одинаковые \code{grid}, \code{shape}, \code{m} и \code{dpl/kd}. Если экспорт
сделан под одну задачу, а ADDA запущена на другой, результат нельзя считать
валидным.

\chapter{Предобуславливание в ADDA}

\section{Левое предобуславливание}

В патче ADDA для \code{BiCGStab} реализовано левое предобуславливание:

\begin{equation}
  M A x = M b.
\end{equation}

Это значит, что в итерационном методе многие векторы живут в
предобусловленном пространстве. В \file{adda_src_modified/iterative.c} есть
отдельная логика для сохранения фактической невязки исходной системы, потому
что внутренняя невязка $M r$ не равна физической невязке $r=b-Ax$.

\section{Почему предобуславливатель может замедлять}

Пусть baseline требует $k_0$ итераций, а run с предобуславливателем требует
$k_M$. Пусть стоимость одной baseline-итерации равна $C_A$, а стоимость
итерации с предобуславливателем равна $C_A+C_M$. Тогда грубая модель:

\begin{equation}
  T_0 \approx k_0 C_A,
\end{equation}

\begin{equation}
  T_M \approx T_{\mathrm{import}} + k_M(C_A+C_M).
\end{equation}

Ускорение есть только если:

\begin{equation}
  k_M(C_A+C_M) + T_{\mathrm{import}} < k_0 C_A.
\end{equation}

Поэтому недостаточно уменьшить число итераций на 20\%, если каждая итерация
стала в 2 раза дороже.

\section{Wall time, iterations и residual}

Нужно различать три метрики:

\begin{description}[leftmargin=3.5cm,style=nextline]
  \item[Iterations]
  Сколько итераций сделал решатель. Полезно, но не равно времени.

  \item[Final residual]
  К какой невязке пришёл решатель. Если оба run упёрлись в \code{maxiter},
  final residual может быть важнее формального speedup.

  \item[Wall time]
  Реальное время по часам. Главная метрика для пользователя.
\end{description}

Если видим:

\begin{lstlisting}
hex_DL1_g48: 300->300 (1.0x)
\end{lstlisting}

это означает, что baseline и preconditioned run оба упёрлись в лимит 300
итераций. Такой чекпоинт не должен становиться \code{best_model.pt}.

\chapter{Архитектура ConvSAI/K2}

\section{Назначение}

ConvSAI/K2 --- универсальная пространственная модель. Она по форме частицы и
физическим параметрам предсказывает локальное свёрточное ядро
предобуславливателя. Главный класс:

\begin{lstlisting}
neural_precond.model.ConvSAI_Universal
\end{lstlisting}

Основной checkpoint:

\begin{lstlisting}
models/k2v3/checkpoints/best_model.pt
\end{lstlisting}

\section{Входы модели}

Модель получает:

\begin{itemize}[leftmargin=2em]
  \item occupancy grid формы: бинарный тензор $0/1$;
  \item $m_{\mathrm{re}}$;
  \item $m_{\mathrm{im}}$;
  \item $kd$;
  \item $\log(\mathrm{grid})$.
\end{itemize}

Функция \code{positions_to_occupancy} переводит позиции диполей в 3D occupancy:

\begin{equation}
  \mathrm{positions} \rightarrow X \in \{0,1\}^{1\times1\times G\times G\times G}.
\end{equation}

\section{ShapeEncoder3D}

Геометрия кодируется небольшой 3D CNN:

\begin{lstlisting}
Conv3d(1 -> 16), ReLU, MaxPool3d
Conv3d(16 -> 32), ReLU, MaxPool3d
Conv3d(32 -> 64), ReLU, MaxPool3d
AdaptiveAvgPool3d(1)
Linear(64 -> shape_embed_dim)
\end{lstlisting}

Перед encoder occupancy приводится к фиксированному разрешению
\code{encoder_resolution=32}. Это позволяет подавать разные grid, но означает,
что очень тонкие геометрические детали могут сглаживаться.

\section{MLP и stencil}

После encoder модель склеивает:

\begin{equation}
  [m_{\mathrm{re}}, m_{\mathrm{im}}, kd, \log(g), s],
\end{equation}

где $s$ --- shape embedding. При \code{shape_embed_dim=16} вход MLP имеет
20 чисел.

MLP выдаёт комплексное ядро:

\begin{equation}
  K[\Delta x,\Delta y,\Delta z] \in \C^{3\times3}.
\end{equation}

При \code{r_cut=7} stencil содержит 1419 смещений. На каждом смещении хранится
3 на 3 комплексный блок.

\section{Что значит K2}

K2 или K$^2$ означает, что модель предсказывает ядро $K$, но применяется:

\begin{equation}
  \widehat{M}(k) = \widehat{K}(k)\,\widehat{K}(k).
\end{equation}

В пространстве это увеличивает эффективный радиус примерно с 7 до 14 клеток.
Параметров модели не становится больше, но экспортируемый stencil после
возведения в квадрат становится шире.

\section{Сильные и слабые стороны}

\begin{longtable}{p{0.3\textwidth}p{0.6\textwidth}}
\toprule
Сторона & Практический смысл \\
\midrule
Плюс & Один checkpoint можно использовать для разных форм и параметров. \\
Плюс & Mode 3 файл обычно компактнее spectral-export с большим radius. \\
Плюс & Простая архитектура и быстрый forward. \\
Минус & Effective radius около 14 может быть недостаточен для больших grid. \\
Минус & Универсальность снижает специализацию под hard hex-prism cases. \\
Минус & Loss может улучшаться без улучшения real ADDA wall time. \\
\bottomrule
\end{longtable}

\chapter{Архитектура Spectral ConvSAI}

\section{Главная идея}

Spectral ConvSAI работает не как локальный spatial stencil predictor, а как
частотный predictor. Для каждой точки FFT-сетки модель предсказывает:

\begin{equation}
  \widehat{M}(k) \in \C^{3\times3}.
\end{equation}

Главный класс:

\begin{lstlisting}
neural_precond.model.ConvSAI_Spectral
\end{lstlisting}

Ветка:

\begin{lstlisting}
spectral-convsai
\end{lstlisting}

\section{Вход на каждой частоте}

Для каждой частоты модель получает:

\begin{itemize}[leftmargin=2em]
  \item полный комплексный $D_{\mathrm{hat}}(k)$, разложенный в 18 вещественных чисел;
  \item координаты частоты;
  \item глобальный conditioning-вектор от формы и физических параметров.
\end{itemize}

Глобальный conditioning строится так же концептуально: shape encoder плюс
$m_{\mathrm{re}}, m_{\mathrm{im}}, kd, \log(g)$.

\section{Почему spectral лучше масштабируется}

ConvSAI/K2 ограничен spatial radius. Spectral-модель сразу работает в
частотной области, где operator inverse естественно описывается как функция
частоты. Поэтому на больших grid и сложных refractive indices spectral-модель
часто сильнее.

Цена:

\begin{itemize}[leftmargin=2em]
  \item экспорт зависит от конкретной задачи;
  \item файл может быть больше;
  \item import и apply могут стать существенной частью wall time.
\end{itemize}

\section{Transformer в spectral-модели}

Transformer здесь не является моделью, которая напрямую выдаёт весь
предобуславливатель. Он используется как coarse frequency context.

Если включить:

\begin{lstlisting}
--spectral_transformer_tokens 8
\end{lstlisting}

то берётся грубая сетка частот:

\begin{equation}
  i_x,i_y,i_z = \mathrm{linspace}(0, g-1, n_{\mathrm{tokens}}).
\end{equation}

Для выбранных точек собираются признаки $D_{\mathrm{hat}}$, координаты частоты
и квадратичные признаки. Transformer выдаёт поправку к global conditioning,
после чего обычный per-frequency MLP предсказывает $\widehat{M}(k)$ на всех
частотах.

\section{Correction branch}

В spectral-модели есть zero-initialized additive correction branch. Она нужна,
чтобы дообучать небольшую поправку, не разрушая базовый checkpoint. Практически
это безопаснее, чем размораживать всю модель, когда уже есть рабочая spectral
версия.

\chapter{Экспорт в ADDA}

\section{Почему нужен export}

ADDA не исполняет PyTorch. Она читает бинарный файл \code{.precond}. Поэтому
любая модель проходит цепочку:

\begin{equation}
  \text{checkpoint} \rightarrow \text{export script} \rightarrow \text{.precond} \rightarrow \text{ADDA}.
\end{equation}

Если export script не соответствует архитектуре checkpoint, результат либо не
загрузится, либо будет математически неправильным.

\section{ConvSAI/K2 export}

\begin{lstlisting}
python3 apps/export_universal_precond.py \
  --checkpoint models/k2v3/checkpoints/best_model.pt \
  --squared_kernel \
  --shape prism \
  --ay 6.0 \
  --az 1.0 \
  --grid 80 \
  --m_re 2.5 \
  --m_im 0.0 \
  --kd 0.41887902047863906 \
  --output exports/prism_g80_m25_k2v3.precond
\end{lstlisting}

Здесь \code{--squared_kernel} обязателен для K2 checkpoint, иначе экспорт будет
соответствовать другой математике.

\section{Spectral export}

\begin{lstlisting}
python3 apps/export_spectral_precond.py \
  --checkpoint models/spectral/checkpoints/best_hex_prism_real_32to96_r40_quality10h_sym.pt \
  --shape prism \
  --ay 6.0 \
  --az 1.0 \
  --grid 80 \
  --m_re 2.5 \
  --m_im 0.0 \
  --kd 0.41887902047863906 \
  --threshold-rel 1e-6 \
  --max-radius 40 \
  --blend-identity 1.0 \
  --symmetry z180_zflip \
  --output exports/prism_g80_m25_spectral.precond
\end{lstlisting}

Параметры \code{--threshold-rel}, \code{--max-radius} и \code{--symmetry}
контролируют, сколько spatial-компонент будет сохранено после преобразования
из частотного представления.

\section{Правило совпадения параметров}

\begin{longtable}{p{0.35\textwidth}p{0.45\textwidth}}
\toprule
Export & ADDA \\
\midrule
\code{--grid 80} & \code{-grid 80} \\
\code{--shape prism --ay 6 --az 1} & \code{-shape prism 6.0 1.0} \\
\code{--m_re 2.5 --m_im 0.0} & \code{-m 2.5 0.0} \\
\code{--kd 0.41887902047863906} & \code{-dpl 15} \\
\bottomrule
\end{longtable}

Несовпадение этих параметров является одной из самых частых причин
``ускорения без ускорения''.

\chapter{Форматы precond-файлов}

\section{Поддерживаемые modes}

В \file{adda_src_modified/precond.h} определены основные modes:

\begin{longtable}{p{0.15\textwidth}p{0.7\textwidth}}
\toprule
Mode & Что хранит \\
\midrule
0 & ILU factors. \\
1 & Sparse approximate inverse matrix. \\
2 & Polynomial preconditioner. \\
3 & ConvSAI spatial stencil. \\
4 & FFTDIRECT: готовый frequency-domain kernel. \\
5 & FFTDIRECT\_XSLAB\_F32: x-slab float32 layout для MPI. \\
\bottomrule
\end{longtable}

\section{Mode 3}

Mode 3 хранит spatial stencil:

\begin{equation}
  \{(\Delta x_s,\Delta y_s,\Delta z_s), K_s\}_{s=1}^{S}.
\end{equation}

При импорте ADDA строит frequency-domain kernel через FFT. Это компактно, но
на больших stencil и MPI import может быть заметным.

\section{Mode 4}

Mode 4 хранит уже готовый $\widehat{M}$ на FFT-сетке. Import быстрее, потому
что ADDA не надо пересобирать $\widehat{M}$ из spatial stencil.

\begin{lstlisting}
python3 apps/convert_convsai_to_fftdirect.py \
  --input exports/prism_g80_m25_spectral.precond \
  --output exports/prism_g80_m25_spectral.fftdirect.precond \
  --grid-x 160 \
  --grid-y 192 \
  --grid-z 160
\end{lstlisting}

Важно: \code{grid-x/y/z} --- это FFT grid из ADDA log, а не просто particle
\code{-grid}.

\section{Mode 5}

Mode 5 хранит frequency kernel в float32 по x-slab. Это полезно для MPI, потому
что каждый процесс может читать свою часть, а не весь dense kernel.

\begin{lstlisting}
python3 apps/convert_fftdirect_to_xslab_f32.py \
  --input exports/prism_g80_m25_spectral.fftdirect.precond \
  --output exports/prism_g80_m25_spectral.xslab_f32.precond
\end{lstlisting}

Mode 4/5 ускоряют import и уменьшают overhead. Они не исправляют плохой
математический предобуславливатель.

\chapter{Сборка и запуск ADDA}

\section{Сборка patched ADDA}

Репозиторий не хранит весь upstream ADDA. Нужно взять чистую ADDA и наложить
патч:

\begin{lstlisting}
git clone https://github.com/adda-team/adda adda
cp adda_src_modified/* adda/src/

make -C adda/src seq mpi \
  FFTW3_INC_PATH="$HOME/.local/include" \
  FFTW3_LIB_PATH="$HOME/.local/lib"
\end{lstlisting}

\section{Sequential run}

\begin{lstlisting}
export LD_LIBRARY_PATH="$HOME/.local/lib:${LD_LIBRARY_PATH:-}"

adda/src/seq/adda \
  -dir runs/prism_g80_m25_seq \
  -grid 80 \
  -m 2.5 0.0 \
  -shape prism 6.0 1.0 \
  -dpl 15 \
  -eps 3 \
  -iter bicgstab \
  -precond exports/prism_g80_m25_spectral.precond
\end{lstlisting}

\section{MPI run}

\begin{lstlisting}
export LD_LIBRARY_PATH="$HOME/.local/lib:${LD_LIBRARY_PATH:-}"

mpirun -np 16 adda/src/mpi/adda_mpi \
  -dir runs/prism_g80_m25_mpi \
  -grid 80 \
  -m 2.5 0.0 \
  -shape prism 6.0 1.0 \
  -dpl 15 \
  -eps 3 \
  -iter bicgstab \
  -precond exports/prism_g80_m25_spectral.precond
\end{lstlisting}

\section{FGMRES}

Для больших FFTDIRECT cases иногда полезен FGMRES:

\begin{lstlisting}
ADDA_FGMRES_RESTART=100 \
mpirun -np 16 adda/src/mpi/adda_mpi \
  -dir runs/prism_g80_m25_mpi_fgmres \
  -grid 80 \
  -m 2.5 0.0 \
  -shape prism 6.0 1.0 \
  -dpl 15 \
  -eps 3 \
  -iter fgmres \
  -precond exports/prism_g80_m25_spectral.precond
\end{lstlisting}

FGMRES использует правое предобуславливание и может требовать меньше
preconditioner applications на эффективный шаг, но его надо сравнивать
эмпирически на целевых cases.

\chapter{Обучение}

\section{Файл trainer}

Главный trainer:

\begin{lstlisting}
train_v7/train.py
\end{lstlisting}

Он делает цикл:

\begin{enumerate}[leftmargin=2em]
  \item сэмплирует форму, grid, $m$, $kd$;
  \item строит позиции диполей;
  \item строит \code{FFTMatVec};
  \item строит occupancy grid;
  \item получает kernel или spectral conditioning от модели;
  \item считает loss;
  \item делает backward и AdamW step;
  \item периодически запускает validation.
\end{enumerate}

Базовый seed в trainer:

\begin{lstlisting}
seed = 42
\end{lstlisting}

\section{Точность вычислений}

По умолчанию веса модели обучаются в fp32. Однако многие операции с системой,
FFTMatVec и комплексными векторами выполняются в \code{complex128}. Export
пишет комплексные double значения в \code{.precond}. AMP/fp16 в стандартной
команде обучения не включены.

\section{Loss-функции}

\subsection{Probe loss}

Идея:

\begin{equation}
  \mathcal{L}_{\mathrm{probe}} =
  \mathbb{E}_z
  \frac{\|M A z - z\|^2}{\|z\|^2}.
\end{equation}

Это проверяет, похож ли $M$ на левый inverse. Минус: случайные векторы могут
плохо коррелировать с реальными plane-wave RHS в ADDA.

\subsection{Adversarial probe loss}

Ищет плохие направления для $I-MA$ через power iteration, затем учит модель на
них. Это сильнее обычного probe loss, но всё равно не гарантирует real ADDA
speedup.

\subsection{Planewave Krylov loss}

Использует RHS, похожие на ADDA plane wave. Вместо случайной правой части
строятся две поляризации падающей волны. Это ближе к реальной задаче.

\subsection{Planewave BiCGStab loss}

Unrolled left-preconditioned BiCGStab loss. Он пытается оптимизировать именно
то, что делает ADDA solver. Но у него есть опасность: если loss слишком
свободный, модель может улучшать краткий unroll, не оставаясь хорошим
приближением $A^{-1}$.

\subsection{Зачем anchor-термы}

Anchor-термы возвращают модель к смыслу inverse:

\begin{equation}
  M A z \approx z,
\end{equation}

\begin{equation}
  A M z \approx z.
\end{equation}

Практический вывод: для ConvSAI/K2 fine-tune на \code{planewave_bicgstab}
опасно ставить:

\begin{lstlisting}
--anchor_probe_weight 0.0
--anchor_right_probe_weight 0.0
\end{lstlisting}

Такой запуск может показывать улучшающийся loss, но не давать ускорения в
ADDA. Более разумный старт:

\begin{lstlisting}
--anchor_probe_weight 0.1
--anchor_right_probe_weight 0.1
--anchor_num_probes 1
--anchor_probe_chunk 1
\end{lstlisting}

\chapter{Guarded validation}

\section{Зачем gate}

Нельзя сохранять \code{best_model.pt} только потому, что training loss стал
лучше. Нужна схема:

\begin{equation}
  \text{train chunk}
  \rightarrow
  \text{export candidate}
  \rightarrow
  \text{ADDA MPI validation}
  \rightarrow
  \text{promote only if better}.
\end{equation}

Именно для этого используется:

\begin{lstlisting}
--guarded_val
--solve_val_interval 250
--guarded_min_large_speedup 1.01
\end{lstlisting}

\section{Исправленная проблема ConvSAI/K2 validation}

В текущей рабочей копии исправлена важная ошибка: large-grid ADDA validation
раньше был жёстко завязан на spectral exporter. Для ConvSAI/K2 это неправильно.

Правильная логика:

\begin{lstlisting}
if spectral:
    apps/export_spectral_precond.py
else:
    apps/export_universal_precond.py
\end{lstlisting}

Также для \code{SquaredConvSAI} validation checkpoint сохраняется без префикса
\code{base.}, чтобы \code{export_universal_precond.py} мог его загрузить.

\section{Фильтрация validation по форме}

Если обучение запущено так:

\begin{lstlisting}
--only_shape hex_prism
--hex_dl_min 1.0
--hex_dl_max 1.0
\end{lstlisting}

то validation должен проверять hex-prism DL=1, а не sphere. Поэтому добавлена
фильтрация fixed validation configs по \code{only_shape} и диапазону DL.

\section{Как читать gate}

Пример плохого, но корректно обработанного результата:

\begin{lstlisting}
Validation (3 configs) speedup: 1.12x
hex_DL1_g32: 1000->751 (1.3x) | hex_DL1_g40: 1000->1000 (1.0x) | hex_DL1_g48: 300->300 (1.0x)
Guarded best gate: FAIL; required large-grid ADDA cases did not pass: hex_DL1_g48
>>> Best update skipped: guarded large-grid gate failed
\end{lstlisting}

Это значит, что trainer не испортил \code{best_model.pt}. Periodic
\code{model_step500.pt} может существовать, но он не является хорошей моделью.

\chapter{Практические команды обучения}

\section{ConvSAI/K2 baseline training}

\begin{lstlisting}
python3 train_v7/train.py \
  --name convsai_universal_k2_v3 \
  --device 0 \
  --save \
  --loss adversarial \
  --squared_kernel \
  --r_cut 7 \
  --hidden_size 512 \
  --num_layers 4 \
  --grid_min 8 \
  --grid_max 64 \
  --m_re_min 1.5 \
  --m_re_max 4.0 \
  --kd_min 0.2 \
  --kd_max 0.8 \
  --ema_decay 0.999 \
  --warmup_steps 500 \
  --num_steps 60000 \
  --lr 5e-4
\end{lstlisting}

\section{ConvSAI/K2 large-grid gated fine-tune with anchors}

\begin{lstlisting}
PYTORCH_CUDA_ALLOC_CONF=expandable_segments:True \
python3 -u train_v7/train.py \
  --resume models/k2v3/checkpoints/best_model.pt \
  --name convsai_k2_hex_g32to96_adda_gated_anchor \
  --device 0 \
  --save \
  --loss planewave_bicgstab \
  --squared_kernel \
  --r_cut 7 \
  --hidden_size 512 \
  --num_layers 4 \
  --only_shape hex_prism \
  --hex_dl_min 1.0 \
  --hex_dl_max 1.0 \
  --grid_min 32 \
  --grid_max 96 \
  --m_re_min 2.0 \
  --m_re_max 3.2 \
  --m_im_min 0.0 \
  --m_im_max 0.0 \
  --kd_min 0.41887902047863906 \
  --kd_max 0.41887902047863906 \
  --krylov_iters 1 \
  --anchor_probe_weight 0.1 \
  --anchor_right_probe_weight 0.1 \
  --anchor_num_probes 1 \
  --anchor_probe_chunk 1 \
  --lr 2e-6 \
  --warmup_steps 100 \
  --ema_decay 0.0 \
  --curriculum_frac 0.0 \
  --guarded_val \
  --guarded_min_large_speedup 1.01 \
  --solve_val_interval 250 \
  --val_interval 0 \
  --adda_val_np 16 \
  --adda_val_eps 3 \
  --adda_val_maxiter_base 300 \
  --adda_val_maxiter_precond 300 \
  --adda_val_timeout 600 \
  --adda_mpi_bin adda/src/mpi/adda_mpi \
  --fftw_lib_path "$HOME/.local/lib" \
  --skip_final_eval \
  --log_interval 10 \
  --save_interval 250 \
  --num_steps 5000
\end{lstlisting}

\section{Spectral continuation loop}

В ветке \code{spectral-convsai} используется promotion loop:

\begin{lstlisting}
DEVICE=0 \
SEED_BASE=6200 \
CYCLES=200 \
STEPS_PER_CYCLE=25 \
LR=2e-6 \
TRAIN_GRID_MIN=32 \
TRAIN_GRID_MAX=96 \
VAL_GRIDS=80,96 \
SCORE_MODE=max \
M_RE=3.0 \
M_IM=0.0 \
DPL=15 \
KD=0.41887902047863906 \
RADIUS=40 \
LOSS=planewave_bicgstab \
KRYLOV_ITERS=1 \
ANCHOR_PROBE_WEIGHT=0.0 \
ANCHOR_RIGHT_PROBE_WEIGHT=0.0 \
SPECTRAL_FREQ_CHUNK_SIZE=65536 \
SPECTRAL_FREQ_CHECKPOINT_CHUNKS=1 \
CURRICULUM_FRAC=0.0 \
NP=16 \
VAL_MAXITER=120 \
VAL_TIMEOUT=600 \
EXPORT_SYMMETRY=z180_zflip \
START_CHECKPOINT=models/spectral/checkpoints/best_hex_prism_real_32to96_r40_quality10h_sym.pt \
BEST_CHECKPOINT=models/spectral/checkpoints/best_hex_prism_real_32to96_r40_quality10h_sym_resume.pt \
BEST_SCORE_FILE=models/spectral/checkpoints/best_hex_prism_real_32to96_r40_quality10h_sym_resume.score \
NAME_PREFIX=SPECTRAL_RESUME_G32TO96_SYM_R40 \
RUN_ROOT=runs/SPECTRAL_RESUME_G32TO96_SYM_R40 \
./train_spectral_hex_real_loop.sh
\end{lstlisting}

\chapter{Как смотреть логи}

\section{Процессы}

\begin{lstlisting}
ps -eo pid,ppid,stat,etime,cmd | rg 'train_v7|adda_mpi|mpirun'
\end{lstlisting}

\section{GPU}

\begin{lstlisting}
nvidia-smi
\end{lstlisting}

\section{Лог обучения}

\begin{lstlisting}
tail -f logs/<RUN_NAME>/train.log
\end{lstlisting}

\section{Файлы результата}

\begin{lstlisting}
find results/<RUN_NAME> -maxdepth 1 -type f -printf '%f %s\n' | sort
\end{lstlisting}

\section{Типичная строка training log}

\begin{lstlisting}
[500/5000] loss=-1.205630 gnorm=10.9315 lr=2.0e-06 1.4 steps/s (hex_DL1.00 g48 m=2.70+0.00i kd=0.42)
\end{lstlisting}

\begin{description}[leftmargin=3.2cm,style=nextline]
  \item[\code{loss}]
  Внутренний training loss. Не равен speedup.

  \item[\code{gnorm}]
  Норма градиента после clipping.

  \item[\code{steps/s}]
  Скорость обучения. На больших grid она резко ниже.

  \item[\code{hex_DL1.00 g48}]
  Текущая сэмплированная задача.
\end{description}

\chapter{Heatmap и benchmark}

\section{Что должна показывать heatmap}

Хорошая heatmap должна измерять не только факт запуска, а минимум:

\begin{itemize}[leftmargin=2em]
  \item baseline wall time;
  \item preconditioned wall time;
  \item speedup;
  \item iterations;
  \item final residual;
  \item converged/capped status;
  \item export/import overhead, если это часть workflow.
\end{itemize}

\section{Три разных speedup}

\begin{description}[leftmargin=4.2cm,style=nextline]
  \item[\code{heatmap_adda_wall_speedup}]
  Ускорение самого ADDA solve.

  \item[\code{heatmap_precond_wall_s}]
  Абсолютное время preconditioned run.

  \item[\code{heatmap_total_elapsed_speedup}]
  Ускорение полного pipeline, если туда входят export/import.
\end{description}

Для одного запуска total elapsed speedup может быть плохим, потому что export
и import не успевают окупиться. Для \code{-orient avg} на сотни ориентаций
import/export amortized, и картина может быть лучше.

\section{Ориентации}

Ускорение может зависеть от orientation, потому что меняется правая часть,
путь сходимости и взаимодействие с симметрией частицы. Для честной оценки
\code{-orient avg} нужно тестировать не один угол, а набор ориентаций:

\begin{equation}
  S = \frac{1}{N}\sum_{i=1}^{N}
  \frac{T_{0,i}}{T_{M,i}}.
\end{equation}

Также стоит смотреть worst case:

\begin{equation}
  S_{\min} = \min_i \frac{T_{0,i}}{T_{M,i}}.
\end{equation}

\chapter{Текущие результаты и ограничения}

\section{Spectral}

По актуальным заметкам ветки \code{spectral-convsai}:

\begin{lstlisting}
mean ADDA wall-time speedup: 9.93x
max ADDA wall-time speedup: 23.33x
mean total elapsed speedup including export: 4.43x
max total elapsed speedup including export: 15.87x
\end{lstlisting}

После MPI distributed fix были локальные проверки:

\begin{longtable}{p{0.3\textwidth}p{0.18\textwidth}p{0.18\textwidth}p{0.18\textwidth}}
\toprule
Case & Baseline & New precond & Speedup \\
\midrule
sphere g32, m=3 & 15.31 s & 2.13 s & 7.19x \\
sphere g48, m=3 & 108.26 s & 8.09 s & 13.38x \\
sphere g64, m=2 & 141.24 s & 53.69 s & 2.63x \\
\bottomrule
\end{longtable}

Вывод: 20x возможно в отдельных областях heatmap, но не является постоянной
гарантией по всем grid, m, shape и orientation.

\section{ConvSAI/K2}

Исторически K2 v3 давал большие итерационные ускорения на умеренных задачах,
например около 30x по итерациям на некоторых g33 cases. Но свежий hard
hex-prism test показал плохую картину:

\begin{longtable}{p{0.22\textwidth}p{0.16\textwidth}p{0.18\textwidth}p{0.22\textwidth}}
\toprule
Run & Wall time & Iters & Status \\
\midrule
baseline & 21.291 s & 600 & cap \\
k2 original & 43.588 s & 600 & cap, worse residual \\
k2 step1500 & 42.843 s & 600 & cap, worse residual \\
k2 step2000 & 43.496 s & 600 & cap, worse residual \\
spectral & 38.101 s & 542 & partial convergence \\
\bottomrule
\end{longtable}

Это не значит, что ConvSAI/K2 бесполезен вообще. Это значит, что на hard
large-grid prism case текущая K2 модель не даёт нужного ускорения и не должна
дообучаться вслепую.

\chapter{Почему последние попытки были подозрительными}

\section{Проблема exporter в validation}

Если ConvSAI/K2 validation использует spectral exporter, то real ADDA gate не
проверяет настоящую ConvSAI/K2 модель. Это исправлено в \file{train_v7/train.py}.

\section{Проблема нулевых anchors}

Запуск с:

\begin{lstlisting}
--anchor_probe_weight 0.0
--anchor_right_probe_weight 0.0
\end{lstlisting}

слишком свободен. Он может оптимизировать короткий unrolled BiCGStab path, но
не удерживать $M$ как приблизительный inverse. Симптом совпал с наблюдением:
loss улучшается, а large-grid ADDA gate остаётся \code{300->300}.

\section{Что значит ``чинить'' в этом месте}

Практически правильный фикс:

\begin{enumerate}[leftmargin=2em]
  \item не продолжать blind run;
  \item запускать с anchor weights;
  \item оставить \code{guarded_val};
  \item сохранять \code{best_model.pt} только после ADDA PASS;
  \item если на 750--1000 шагах gate стабильно \code{300->300}, менять подход,
  а не просто ждать 10 часов.
\end{enumerate}

\chapter{Ограничения и направления развития}

\section{Low-rank}

Low-rank блок добавляет несколько глобальных мод:

\begin{equation}
  M \approx M_{\mathrm{local}} + U C V^*,
\end{equation}

где локальная часть ловит ближнюю физику, а low-rank часть ловит дальние
медленные modes. Это естественное развитие для ConvSAI/K2, потому что его
главная слабость --- ограниченный radius.

\section{Deflation}

Deflation явно убирает из Krylov-решателя плохие медленно сходящиеся
направления. Идея:

\begin{equation}
  \mathcal{K} \rightarrow \mathcal{K} \oplus \mathrm{span}(u_1,\dots,u_r).
\end{equation}

Это может быть полезно, если heatmap показывает повторяющиеся bad modes.

\section{SO(3) equivariance}

Обычный 3D CNN не гарантирует корректного поведения при произвольных
поворотах. Для \code{-orient avg} полезна архитектура или аугментация, которая
учитывает вращения:

\begin{equation}
  f(Rx) = R f(x) R^{-1}.
\end{equation}

Для призм также важны дискретные симметрии, например \code{z180} и
\code{zflip}.

\section{Analytic circulant inverse}

Так как FFT-оператор уже имеет block-circulant структуру, можно строить часть
предобуславливателя аналитически в частотной области, а нейросетью учить
только поправку. Spectral branch ближе всего к этой идее.

\section{Аугментации}

Нужные аугментации для ориентаций:

\begin{itemize}[leftmargin=2em]
  \item случайные \code{alpha beta gamma};
  \item повороты occupancy grid;
  \item symmetry-aware export для призмы;
  \item validation на нескольких orientation seeds;
  \item отдельный average и worst-case speedup.
\end{itemize}

\chapter{Практический чеклист}

\section{Перед запуском}

\begin{enumerate}[leftmargin=2em]
  \item Убедиться, что ветка правильная: \code{convsai-universal} или
  \code{spectral-convsai}.
  \item Собрать patched ADDA.
  \item Проверить \code{LD_LIBRARY_PATH}.
  \item Проверить, что export параметры совпадают с ADDA параметрами.
  \item Проверить, что используется правильный exporter.
\end{enumerate}

\section{Перед доверием к checkpoint}

\begin{enumerate}[leftmargin=2em]
  \item Это \code{best_model.pt}, а не случайный \code{model_step*.pt}.
  \item В логе есть \code{Guarded best gate: PASS}.
  \item Есть сравнение с baseline на той же задаче.
  \item Проверены wall time, iterations и final residual.
  \item Проверено несколько grid/m/orientation.
\end{enumerate}

\section{Если нет ускорения}

\begin{enumerate}[leftmargin=2em]
  \item Проверить совпадение export/run параметров.
  \item Проверить, не упёрлись ли оба run в \code{maxiter}.
  \item Проверить final residual.
  \item Проверить preconditioner apply time.
  \item Проверить import overhead.
  \item Для ConvSAI/K2 проверить anchors и guarded validation.
  \item Для Spectral проверить radius, threshold и symmetry.
\end{enumerate}

\chapter{Карта файлов}

\begin{longtable}{p{0.38\textwidth}p{0.52\textwidth}}
\toprule
Файл или папка & Назначение \\
\midrule
\file{README.md} & Краткая инструкция текущей ветки. \\
\file{docs/PROJECT_BOOK_RU.tex} & Этот LaTeX-документ. \\
\file{docs/PROJECT_BOOK_RU.md} & Старый Markdown-конспект, оставлен как черновик. \\
\file{train_v7/train.py} & Trainer, validation, checkpoint logic. \\
\file{neural_precond/model.py} & ConvSAI, Spectral, ShapeEncoder, transformer. \\
\file{neural_precond/loss.py} & Loss-функции. \\
\file{apps/export_universal_precond.py} & ConvSAI/K2 export. \\
\file{apps/export_spectral_precond.py} & Spectral export в ветке \code{spectral-convsai}. \\
\file{apps/convert_convsai_to_fftdirect.py} & Mode 3 -> mode 4. \\
\file{apps/convert_fftdirect_to_xslab_f32.py} & Mode 4 -> mode 5. \\
\file{adda_src_modified/precond.c} & Загрузка и применение preconditioner в ADDA. \\
\file{adda_src_modified/iterative.c} & Изменения solver для preconditioning. \\
\file{models/k2v3/checkpoints/best_model.pt} & ConvSAI/K2 checkpoint. \\
\bottomrule
\end{longtable}

\backmatter

\chapter{Итоговая модель понимания}

Если сжать весь проект до десяти пунктов:

\begin{enumerate}[leftmargin=2em]
  \item ADDA решает $Ax=b$ итерационным методом.
  \item Предобуславливатель $M$ должен делать $MA$ ближе к identity.
  \item Реальная метрика --- wall time, а не training loss.
  \item ConvSAI/K2 --- универсальный local spatial stencil predictor.
  \item Spectral --- frequency-domain predictor, лучше подходящий для больших grid.
  \item Export должен точно соответствовать shape/grid/m/kd.
  \item Mode 4/5 ускоряют import, но не чинят плохой $M$.
  \item Checkpoint должен проходить real ADDA gate.
  \item Для \code{-orient avg} нужно тестировать много ориентаций.
  \item Для цели 20x основной кандидат сейчас --- Spectral branch; ConvSAI/K2
  требует осторожного gated fine-tune или архитектурного усиления.
\end{enumerate}

\end{document}
